\documentclass[a4paper,11pt]{article}
\usepackage[a4paper,margin=2cm]{geometry}
\usepackage[T1]{fontenc}
\usepackage[utf8]{inputenc}
\usepackage[ngerman,english]{babel}
\usepackage{lmodern}
\usepackage{microtype}
\usepackage{xcolor}
\usepackage{parskip}
\usepackage{enumitem}
\usepackage[most]{tcolorbox}
\usepackage{titlesec}
\usepackage[hidelinks]{hyperref}
\definecolor{HeaderBlueDark}{RGB}{70,110,170}
\definecolor{RowBlueLight}{RGB}{235,243,255}
\definecolor{RowBlueDark}{RGB}{220,233,250}
\definecolor{SoftGray}{RGB}{90,90,90}
\setlength{\parindent}{0pt}
\setlist[itemize]{leftmargin=1.4em,itemsep=0.35em,topsep=0.35em}
\linespread{1.06}
\titleformat{\section}[block]{\Large\bfseries\color{HeaderBlueDark}}{}{0pt}{}
\titlespacing{\section}{0pt}{1.3em}{0.8em}
\titleformat{\subsection}[block]{\large\bfseries\color{HeaderBlueDark}}{}{0pt}{}
\titlespacing{\subsection}{0pt}{1.0em}{0.6em}
\newtcolorbox{exbox}{enhanced,breakable,colback=RowBlueLight,colframe=RowBlueDark,boxrule=0.4pt,arc=1.5mm,left=1.2mm,right=1.2mm,top=1mm,bottom=1mm,title={Beispiele \textnormal{(Examples)}},fonttitle=\bfseries,colbacktitle=RowBlueDark,coltitle=black}
\NewDocumentEnvironment{grammarentry}{mm+b}{%
\begin{tcolorbox}[enhanced,breakable,colback=white,colframe=HeaderBlueDark,boxrule=0.6pt,arc=1.5mm,left=1.5mm,right=1.5mm,top=1mm,bottom=1mm,colbacktitle=HeaderBlueDark,coltitle=white,fonttitle=\bfseries,title={#1\hfill\normalfont\footnotesize #2}]
#3
\end{tcolorbox}}{}
\newcommand{\GermanText}[1]{\textbf{Deutsch: }#1\par}
\newcommand{\EnglishText}[1]{{\small\color{SoftGray}\textbf{English: }#1\par}}
\newcommand{\ExampleItem}[2]{\item \textbf{#1}\\{\small\color{SoftGray}#2}}


\title{Relativwort: \glqq wieso\grqq{}}

\date{}

\begin{document}

\maketitle

\tableofcontents

\newpage

\section{Grundbedeutung und Kategorie}

\begin{grammarentry}{wieso}{Relativadverb, kausal, eher umgangssprachlich}

\GermanText{\textbf{wieso} kann in relativer Funktion einen Grund verbinden. Seine ungefähre Bedeutung ist \glqq why; for which reason\grqq{}.}

\EnglishText{\textbf{wieso} can in relative use connect a reason. Its approximate meaning is \glqq why; for which reason\grqq{}.}

\end{grammarentry}

\section{Wichtig: Pronomen oder Adverb?}

\begin{grammarentry}{Grammatische Einordnung}{kein deklinierbares Pronomen}

\GermanText{\textbf{wieso} ist in dieser Verwendung kein deklinierbares Relativpronomen. Es ist ein Relativadverb bzw. relatives Verbindungswort. Deshalb hat es keinen Nominativ, Akkusativ, Dativ oder Genitiv.}

\EnglishText{\textbf{wieso} is not a declinable relative pronoun in this use. It is a relative adverb or relative linking word. Therefore, it has no nominative, accusative, dative, or genitive case.}

\GermanText{Für formelles B1-Schreiben sind \textbf{warum} oder \textbf{weshalb} meistens die sichereren Formen.}

\EnglishText{For formal B1 writing, \textbf{warum} or \textbf{weshalb} are usually the safer forms.}

\end{grammarentry}

\section{Beispiele}

\begin{exbox}

\begin{itemize}

\ExampleItem{Ich kenne den Grund, wieso er gegangen ist.}{I know the reason why he left.}

\ExampleItem{Das ist ein Grund, wieso ich vorsichtig bin.}{That is one reason why I am careful.}

\end{itemize}

\end{exbox}

\section{Wortstellung und Komma}

\begin{grammarentry}{Nebensatzregel}{Verb am Ende}

\GermanText{Der eingeleitete Satz ist ein Nebensatz. Das konjugierte Verb steht normalerweise am Ende, und der Satz wird durch Kommas abgetrennt.}

\EnglishText{The introduced clause is a subordinate clause. The conjugated verb normally stands at the end, and the clause is separated by commas.}

\end{grammarentry}

\section{Merksatz}

\begin{grammarentry}{wieso}{kurz und wichtig}

\GermanText{\textbf{wieso} kann relativ verwendet werden, ist aber grammatisch ein \textbf{Relativadverb/Relativwort}, kein Relativpronomen wie \textbf{wer} oder \textbf{was}.}

\EnglishText{\textbf{wieso} can be used relatively, but grammatically it is a \textbf{relative adverb/relative word}, not a relative pronoun like \textbf{wer} or \textbf{was}.}

\end{grammarentry}

\end{document}
